% Compile: xelatex
\documentclass[10pt]{article}
\usepackage[a4paper, left=2.5cm, right=2.5cm, top=2cm, bottom=2cm]{geometry}
\usepackage{fontspec}
\setmainfont{Calibri}
\usepackage{array}
\usepackage{tabularx}
\usepackage{amssymb}
\usepackage{setspace}
\setstretch{1.1}
\pagenumbering{gobble}
\setlength{\parindent}{0pt}
\setlength{\parskip}{3pt}

\newcommand{\dotline}[1]{\makebox[#1]{\dotfill}}
\newcolumntype{C}[1]{>{\centering\arraybackslash}p{#1}}

\begin{document}

% ===== JUDUL =====
\begin{center}
{\fontsize{13}{15}\selectfont\bfseries
PERNYATAAN TRANSPARANSI\\
PENGGUNAAN KECERDASAN ARTIFISIAL\\
GEMASTIK XIX TAHUN 2026}
\end{center}

\vspace{0.8em}

% ===== IDENTITAS =====
Yang bertandatangan di bawah ini:

\vspace{0.3em}

\begin{tabularx}{\textwidth}{@{}p{4.2cm}l@{}}
Nama Ketua Tim & : \dotline{9.5cm} \\[3pt]
Nomor Induk Mahasiswa & : \dotline{9.5cm} \\[3pt]
Nama Tim & : \dotline{9.5cm} \\[3pt]
Program Studi & : \dotline{9.5cm} \\[3pt]
Nama Dosen Pembimbing & : \dotline{9.5cm} \\[3pt]
Perugruan Tinggi & : \dotline{9.5cm} \\
\end{tabularx}

\vspace{0.5em}

Dengan ini menyatakan bahwa karya tim kami, yang diikutsertakan dalam GEMASTIK XIX Tahun 2026:

\vspace{0.3em}

\begin{tabularx}{\textwidth}{@{}p{4.2cm}l@{}}
Judul Karya & : \dotline{9.5cm} \\[3pt]
Cabang Kompetisi GEMASTIK & : \dotline{9.5cm} \\
\end{tabularx}

\vspace{0.5em}

% ===== PARAGRAF TRANSPARANSI =====
Dengan ini kami menyampaikan pernyataan transparansi penggunaan Kecerdasan Artifisial (KA)/\textit{Artificial Intelligence} (AI) dan Kecerdasan Artifisial Generatif (KA-Gen)/\textit{Generative Artificial Intelligence} (Gen-AI), sebagai berikut:

\vspace{0.3em}

% ===== CHECKBOX 1: CLOUD AI =====
$\square$~Karya inovasi menggunakan sumber daya KA/AI berbasis awan (\textit{Cloud AI}) atau API pihak ketiga sebagai komponen fungsionalitas operasional perangkat, dengan arsitektur sistem, tingkat integrasi, rekayasa instruksi/\textit{prompting}, dan kustomisasinya adalah murni hasil rancangan mandiri tim kami, meliputi:

\vspace{0.3em}

% ===== TABEL 1: AI/API =====
\begin{tabularx}{\textwidth}{|C{0.6cm}|p{3.2cm}|X|X|}
\hline
\textbf{No.} & \textbf{Nama Aplikasi AI/API} & \textbf{Fungsi yang Didukung dalam Inovasi} & \textbf{Kustomisasi/Pengaturan} \\
\hline
1 & & & \\[6pt]
\hline
2 & & & \\[6pt]
\hline
3 & & & \\[6pt]
\hline
4 & Dst\ldots. & & \\[6pt]
\hline
\end{tabularx}

\vspace{0.5em}

% ===== CHECKBOX 2: NASKAH =====
$\square$~Naskah dokumen yang kami unggah dan/atau karya kreasi kami telah disusun sesuai kaidah akademik dan menjunjung integritas ilmiah serta bukan merupakan hasil \textit{generating} oleh Kecerdasan Artifisial Generatif (Gen-AI) secara sepenuhnya, yaitu dengan menggunakan aplikasi deteksi \dotline{2.5cm} yang memberikan hasil deteksi tingkat penggunaan Gen-AI sebesar \dotline{1.2cm}\%, sesuai bukti \textit{AI writing detection} terlampir, yaitu dengan penjelasan dan argumentasi berikut:

\vspace{0.3em}

% ===== TABEL 2: GEN-AI =====
\begin{tabularx}{\textwidth}{|C{0.6cm}|X|}
\hline
\textbf{No.} & \multicolumn{1}{c|}{\textbf{Penjelasan dan Argumentasi Penggunaan Gen-AI}} \\
\hline
1 & \\[6pt]
\hline
2 & \\[6pt]
\hline
3 & \\[6pt]
\hline
4 & \\[6pt]
\hline
5 & Dst\ldots. \\[6pt]
\hline
\end{tabularx}

\vspace{0.8em}

% ===== PENUTUP =====
Demikian surat pernyataan ini saya buat dengan sebenar-benarnya. Apabila di kemudian hari ditemukan bahwa pernyataan ini tidak benar, saya bersedia menerima sanksi sesuai ketentuan yang berlaku pada kompetisi GEMASTIK XIX Tahun 2026.

\vspace{1em}

% ===== TANDA TANGAN =====
\begin{tabularx}{\textwidth}{@{}X X X@{}}
Mengetahui, & & Kota/Kab., Tanggal-Bulan-Tahun \\[3pt]
Pimpinan Bidang Kemahasiswaan & Dosen Pembimbing & Yang Menyatakan \\[45pt]
Nama Lengkap & Nama Lengkap & Nama Lengkap \\[2pt]
NIP/NUPTK \dotline{2.2cm} & NIP/NUPTK \dotline{2.2cm} & NIM/NPM \dotline{2.2cm} \\
\end{tabularx}

\end{document}
